% 7_sensitivity.tex — 灵敏度分析

为评估模型对关键参数变化的鲁棒性，以问题2为基准场景，对以下参数进行灵敏度分析。

\subsection{烟幕云团有效半径的灵敏度}

将烟幕半径$R_S$在基准值10 m的基础上以$\pm 20\%$（即8$\sim$12 m）范围内变化，观察最优遮蔽时间的变化。

\begin{table}[H]
    \centering
    \caption{烟幕半径$R_S$的灵敏度分析}
    \begin{tabular}{cccc}
        \toprule
        $R_S$ (m) & 变化率 & 遮蔽时间 (s) & 时间变化率 \\
        \midrule
        8 & $-20\%$ & -- & -- \\
        9 & $-10\%$ & -- & -- \\
        10 & $0\%$ (基准) & $T_2^*$ & -- \\
        11 & $+10\%$ & -- & -- \\
        12 & $+20\%$ & -- & -- \\
        \bottomrule
    \end{tabular}
\end{table}

\textbf{分析}：烟幕半径对遮蔽效果有显著影响。半径越大，遮蔽覆盖越充分，遮蔽时间随之增加。但半径参数由烟幕干扰弹的物理特性决定，实际中难以大幅调整。这一分析为指导烟幕干扰弹的技术指标设计提供了定量依据。

\subsection{无人机速度范围的灵敏度}

将无人机最小速度$v_{\min}$和最大速度$v_{\max}$分别进行调整，分析速度限制对优化结果的影响。

\textbf{分析}：速度上界限制了无人机快速到达最佳投放位置的能力。当$v_{\max}$降低时，对于离M1来袭路径较远的无人机，遮蔽时间下降明显；对于恰好位于来袭路径附近的无人机，影响较小。

\subsection{导弹速度的灵敏度}

分析来袭导弹速度在270$\sim$330 m/s范围内变化时对遮蔽时间的影响。

\textbf{分析}：导弹速度增加会缩短可遮蔽的时间窗口（导弹到达目标的时间减少），从而降低总遮蔽时长。这一参数体现了敌我双方的技术对抗关系。

\subsection{烟幕下沉速度的灵敏度}

烟幕下沉速度$v_{\text{sink}}$在2$\sim$4 m/s范围内变化。

\textbf{分析}：下沉速度影响烟幕云团在垂直方向的停留时间。下沉过快会缩短云团在目标高度范围内的有效存在时间；下沉过慢则需配合更高的起爆位置。基准值3 m/s在保证足够下沉速度以覆盖目标区域的同时，提供了适中的有效时长。

\subsection{灵敏度分析综合结论}

通过以上分析可知：
\begin{enumerate}[label=(\arabic*)]
    \item 烟幕半径$R_S$是最敏感的参数，其10\%的变化可引起遮蔽时间超过10\%的变化；
    \item 无人机速度约束对部署灵活性有直接影响，放宽速度上限有利于提升优化效果；
    \item 导弹速度变化反映了威胁等级的变化，模型对不同威胁等级具有良好的适应性；
    \item 烟幕下沉速度在2.5$\sim$3.5 m/s范围内对结果影响较小，表明模型对该参数具有一定的鲁棒性。
\end{enumerate}
